
%%%%%%%%%%%%%%%%%%%%%%%%%%%%%%%%%%%%%%%%%%%%%%%%%%%%%%%%%%%%%%%%%%%%%%%%%%%
%%%%%%%%%%%%%%%%%%%%%%%%%%%%%%%%%%%%%%%%%%%%%%%%%%%%%%%%%%%%%%%%%%%%%%%%%%%
%%%%%%%%%%%%%%%%%%%%%%%%%%%%%%%%%%%%%%%%%%%%%%%%%%%%%%%%%%%%%%%%%%%%%%%%%%%

\begin{figure}
    % \begin{minipage}[t]{1.0\textwidth}
    
    %%%%%%%%%%%%%%%%%%%%%%%%%%%%%%%%%%%%%
    %%%%%%%%%%%%%%%%%%%%%%%%%%%%%%%%%%%%%
    
    \begin{algorithm}[H]
        \caption{Sandwich covariance}\algrlabel{sandwich}
        \textbf{Requires: } Differentiable $\theta \mapsto \likhat(\theta)$
        % \textbf{Inputs (global variables):}\\
        % $\theta \mapsto g(\theta)$ \\ %\Comment{Target function}\\
        % $\theta \mapsto \likhat(\theta)$ \\ %\Comment{Log joint distribution}\\
        \begin{algorithmic}
        %
        \State $\thetahat \gets \argmax_{\theta \in \thetadom} \likhat(\theta)$
    %        \Comment{(MAP estimate)}
        \State $\hat{\Sigma}_\theta \gets \infohat^{-1} \scorecovhat \infohat^{-1}$
     %       \Comment{(sandwich estimate of $\cov{\xfdist}{\sqrt{N} \thetahat}$)}
        \State \Return $\gcovmaphat \gets \ggrad{1}(\thetahat)
            \hat{\Sigma}_\theta \ggrad{1}(\thetahat)^\trans$
      %      \Comment{(delta method estimate of
       %         $\cov{\xfdist}{\sqrt{N} \g(\thetahat)}$)}
        % \State \Return $\gcovmaphat$
        %
        \end{algorithmic}
    \end{algorithm}
    %
    % \end{minipage}
    %
    % \hfill\vline\hfill
    %
    % \begin{minipage}[t]{0.33\textwidth}
    %%%%%%%%%%%%%%%%%%%%%%%%%%%%%%%%%%%%%
    %%%%%%%%%%%%%%%%%%%%%%%%%%%%%%%%%%%%%
    
    \begin{algorithm}[H]
        \caption{Bootstrap covariance}\algrlabel{boot}
        % \textbf{Inputs (global variables):}\\
        % $\xvec = (\x_1, \ldots, \x_N)$ \Comment{Data}\\
        % $\theta \mapsto g(\theta)$ \Comment{Target function}\\
        % $M$ \Comment{Number of MCMC samples}\\
        % $B$ \Comment{Number of bootstrap samples}
        \textbf{Requires: } Posterior sampler $\xvec^* \mapsto \theta^1,\ldots,\theta^M \sim \p(\theta | \xvec^*)$
        (e.g. MCMC)
        \begin{algorithmic}
        %
        \For{$b=1,\ldots,B$}
            \State $\xvec^* \gets$ Sample with replacement from the rows of $\xvec$
            \State Sample
                $\theta^1,\ldots,\theta^M \sim \p(\theta | \xvec^*)$
            % \Comment{Note that $\theta^m \in \rdom{D}$.}
            \State $\mathscr{G}^b \gets
                \frac{1}{M} \sum_{m=1}^M g(\theta^m)$
            % \Comment{$\mathscr{G}^b \approx \expect{p(\theta | \xvec^*)}{\g(\theta)}$.
                % Note that $\mathscr{G}^b \in \rdom{\gdim}$.}
        \EndFor
        \State \Return $\gcovboothat \gets \covhat{b \in [B]}{\sqrt{N} \mathscr{G}^b}$
        % \Comment{(Bootstrap estimate of $\cov{\xfdist}{\sqrt{N} \expect{p(\theta | \xvec)}{\g(\theta)}}$)}
        % \State \Return $\gcovboothat$
        %
        \end{algorithmic}
    \end{algorithm}
    %
    %%%%%%%%%%%%%%%%%%%%%%%%%%%%%%%%%%%%%
    %%%%%%%%%%%%%%%%%%%%%%%%%%%%%%%%%%%%%
    
    \begin{algorithm}[H]
        \caption{IJ covariance (our proposal)}\algrlabel{ijalg}
        \textbf{Requires: } Log likelihood $\theta, \x_n \mapsto \ell(\x_n \vert \theta)$,
        posterior samples $\theta^1,\ldots,\theta^M \sim \p(\theta | \xvec)$
        % \textbf{Inputs (global variables):}\\
        % $\xvec = (\x_1, \ldots, \x_N)$ \Comment{Data}\\
        % $\theta^1, \ldots, \theta^M$
        %     \Comment{MCMC samples from $\p(\theta \vert \xvec)$}\\
        % $\theta \mapsto g(\theta)$ \Comment{Target function}\\
        % $x_n, \theta \mapsto \ell(\x_n | \theta)$ \Comment{Log likelihood function}\\
        
        \begin{algorithmic}
        % \algrlabel{IJ Covariance From MCMC}
        %
        \For{$n=1,\ldots,N$}
            \State $\hat{\infl}_n \gets N
                \covhat{m \in [M]}{\ell(\x_n \vert \theta^m), \g(\theta^m)}$
            % \Comment{Approximate posterior covariance}
        \EndFor
        \State \Return $\gcovijhat \gets \covhat{n\in [N]}{\hat{\infl}_n}$
            % \Comment{Approximate covariance across datapoints}
        % \State \Return $\gcovijhat$
        %
        \end{algorithmic}
    \end{algorithm}
    %
        
    %%%%%%%%%%%%%%%%%%%%%%%%%%%%%%%%%%%%%
    %%%%%%%%%%%%%%%%%%%%%%%%%%%%%%%%%%%%%
    % \end{minipage}
    
    \caption{Three algorithms for estimating $\gcovtrue$.}\figlabel{algorithm}
    \end{figure}

There are two well-known classical methods for estimating $\gcovtrue$: the
sandwich covariance and the bootstrap.  The sandwich covariance
requires computation of the maximum a posteriori (MAP) estimate, and is
not available from MCMC samples.  The bootstrap requires running many MCMC
chains with slightly different data sets, which can be computationally
expensive. We describe a third method, the infinitesimal jackknife (IJ)
covariance estimate, which is also classical, but which has not been previously
applied to MCMC estimates of posterior means.  All three
are succinctly summarized in \figref{algorithm}.

% How would you use the Laplace approximation to estimate \gcovtrue?
The sandwich covariance approach to estimating $\gcovtrue$ is given in
\algrref{sandwich}. To compute the sandwich covariance estimate, one first forms
the MAP estimate, $\thetahat$.  Then one estimates the sampling variability of
$\thetahat$ using an asymptotic approximation to the sampling distribution of
M--estimators under misspecification.  From the sampling variability of
$\thetahat$, one then uses the delta method to estimate the sampling variability
of $\g(\thetahat)$ (\citet{stefanski:2002:mestimation}, Chapter 7 of
\citet{van:2000:asymptotic}).\footnote{It would arguably be more standard to
describe the covariance estimator for $\thetahat$ as the ``sandwich covariance''
and the covariance estimator of $\g(\thetahat)$ as the ``delta method covariance.''  But
the IJ covariance estimate is sometimes described as the ``functional delta method''
(\citet{van:2000:asymptotic} Chapter 20), so the authors believe that our
nomenclature is the least ambiguous option.}
% Under the assumption that
% the MAP estimate is able to approximate to the desired posterior expectation,
% $\expect{\post}{\g(\theta)} \approx \g(\thetahat)$, the sampling variability
% of $\g(\thetahat)$ approximates the sampling variability of $\expect{\post}{\g(\theta)}$.
%
When a sufficiently strong BCLT applies, it is known that
$g(\thetahat)$ consistently estimates $\expect{\post}{g(\theta)}$ and so the
sandwich covariance estimate consistently estimates $\gcovtrue$ 
(more details will follow in \secref{finite_dim_clts} below).  However, the MAP
approximation is not parameterization--invariant, and can behave pathologically
in particular datasets, even when there exists a parameterization in which a
BCLT applies.

When the MAP and sandwich covariance are unreliable, one typically turns to MCMC
estimates.  To estimate $\gcovtrue$ from MCMC samples, a classical method is the
bootstrap \citep{efron:1994:bootstrap,huggins:2022:bayesbag}.  The bootstrap is
shown in \algrref{boot}.  The primary difficulty with the bootstrap is that it
requires $B$ different MCMC chains, which can be computationally expensive;
often, in practice, even a single run of MCMC requires considerable computing
resources. When the MCMC noise can be made arbitrarily small, the bootstrap is a
consistent estimator of $\gcovtrue$ in the fixed-dimension case under conditions
given by \citet{huggins:2022:bayesbag}.


% Define the population IJ estimator
We propose an estimator of $\gcovtrue$ that can be
computed from a single MCMC run using only the tractable $\ell(\x_n \vert \z,
\theta)$.  We call our estimator the {\em infinitesimal jackknife (IJ)
covariance estimator}, and denote it as $\gcovij$.  The IJ covariance estimator
is given by the sample variance of a set of simple posterior covariances, as
given in \eqref{ij_terms}.
%
%%%%%%%%%%%%%%%%%%%%%%%%%%%%%%%%%%%%%
%%%%%%%%%%%%%%%%%%%%%%%%%%%%%%%%%%%%%
%\begin{defn}\deflabel{ij_terms}
    %
    % Define the Bayesian empirical influence function and IJ covariance estimate as
    % follows.
    %
    \begin{align}\eqlabel{ij_terms}
        \infl_n := N \cov{\post}{g(\theta), \ell(\x_n \vert \theta)}
        \textrm{ (Influence score)}
        \quad\quad
        \gcovij := \covhat{n \in [N]}{\infl_n}.
        % \textrm{(IJ covariance estimate)}
        % \gcovij &:= \meann(\infl(\x_n) \infl(\x_n)^T - \bar\infl \, \bar\infl^T)
        % &\textrm{(IJ covariance estimate)}\\
        % &\textrm{where }
        % \bar\infl := \meann \infl(\x_n).
    \end{align}
    %
%\end{defn}
%%%%%%%%%%%%%%%%%%%%%%%%%%%%%%%%%%%%%
%%%%%%%%%%%%%%%%%%%%%%%%%%%%%%%%%%%%%
    
% Define the MCMC IJ estimator
In practice, the posterior covariances of \eqref{ij_terms} typically need to be
estimated with MCMC.  We denote  by $\gcovijhat$ the corresponding estimate of $\gcovij$
(see \algrref{ijalg} for a precise statement).
For any fixed $N$, the output $\gcovijhat$ of \algrref{ijalg} approaches
$\gcovij$ under a well-mixing MCMC chain as the number of MCMC samples $M
\rightarrow \infty$. We will theoretically analyze $\gcovij$, though our
experiments will all employ $\gcovijhat$. Importantly, \algrref{ijalg} requires
only a single run of MCMC without relying on accuracy of the MAP, as required by
the sandwich estimator, or running $B$ different MCMC samplers, as required by
the bootstrap estimator. Our key result, given in \thmref{ij_consistent} below,
states general conditions under which $\gcovij \plim \gcovtrue$.
